\newpage\thispagestyle{hl}
\hlbanner{The whole book}{}{Why this book}
\hlsection{Learning from formulation to verified computation}
\begin{hlquote}Computational mechanics is learned well when three things are seen together: how a formulation becomes an algorithm, how the algorithm is implemented, and how its results are verified.\end{hlquote}
\vspace{4pt}\noindent\textbf{From formulation to software.} Finite element analysis is usually taught through weak forms, element matrices, assembly and solution algorithms, while reduced-order modelling is introduced separately through projection, modal truncation or data-driven approximation. The book follows one path from material, mesh and element to equilibrium and post-processing, and on to snapshots, a basis and a reduced prediction, through two packages, \fea{} and \rom{}, that the reader can run and inspect.

\vspace{4pt}\noindent\textbf{Why write the software instead of simply using established tools?} Established finite element and model-reduction libraries serve important purposes, and the book does not claim superiority over them in performance or breadth. It provides a compact implementation that readers can trace, inspect, test and extend, with only Python, NumPy and SciPy needed and PyTorch optional. The two parts are deliberately decoupled: \rom{} works on numerical arrays and does not import \fea{}.

\vspace{4pt}\noindent\textbf{Beyond introductory examples.} The book covers nonlinear equilibrium, contact, transient dynamics, computational acceleration, substructuring, nonlinear reduced models and hyper-reduction for contact, each through its assumptions, implementation choices, evidence and limitations.

\vspace{4pt}\noindent\textbf{Verification and reproducibility.} Verification is part of the method: analytical solutions, independent computations, convergence, consistency checks, automated tests, and errors on held-out cases. Verification is separated from validation, and passing tests are not treated as proof. Every number in the text is produced by a script, so the number a reader sees is the one that was measured, and the book names the methods that are implemented but not exercised by its examples.

\vspace{4pt}\noindent\textbf{A foundation for further work.} The packages are working research software and learning tools, for readers who want to move from using numerical tools to understanding and developing them, not a claim that every method is production-ready.
